\documentclass{article}
\usepackage[utf8]{inputenc}
\usepackage[T1]{fontenc}
\usepackage[spanish]{babel}
\usepackage{graphicx}
\usepackage{float}
\usepackage{array}
\usepackage{booktabs}
\usepackage{geometry}
\usepackage{fancyhdr}
\usepackage{setspace}
\usepackage{parskip}
\usepackage{amsmath}
\usepackage{amsthm} 
\usepackage{amsfonts}
\usepackage{amssymb}
\usepackage{xcolor}
\usepackage{listings}
\usepackage{caption} 
\usepackage{subcaption}
\usepackage{caption}
\usepackage{hyperref}
\usepackage{geometry}
\geometry{margin=1in}
\usepackage{tikz}
\usepackage{eso-pic}
%\usepackage{siunitx}
\usetikzlibrary{positioning, shapes.geometric, arrows.meta}
\graphicspath{ {./Imagenes/} }
\usepackage[pages=some]{background}

\lstset{
    basicstyle=\ttfamily\scriptsize, % Bajamos de small a scriptsize por lo largo de los #pragma
    breaklines=true,                 % Activa el quiebre de línea
    breakatwhitespace=false,         % Cambia a false para que rompa incluso si no hay espacios
    columns=fullflexible,            % Ayuda a que el espaciado sea más natural
    keepspaces=true,                 % Mantiene los espacios del código original
    literate={á}{{\'a}}1 {é}{{\'e}}1 {í}{{\'i}}1 {ó}{{\'o}}1 {ú}{{\'u}}1 {ñ}{{\~n}}1 % Para que acepte acentos en comentarios
}

%\usepackage[backend=biber,style=authoryear]{biblatex}
%\usepackage{csquotes}
%\addbibresource{referencias.bib}

\definecolor{codegreen}{rgb}{0,0.6,0}
\definecolor{codegray}{rgb}{0.5,0.5,0.5}
\definecolor{codeblue}{rgb}{0,0,0.95}
\definecolor{backcolour}{rgb}{0.95,0.95,0.92}

\lstdefinestyle{mystyle}{
    backgroundcolor=\color{backcolour},   
    commentstyle=\color{codegreen},
    keywordstyle=\color{codeblue},
    numberstyle=\tiny\color{codegray},
    stringstyle=\color{codegreen},
    basicstyle=\ttfamily\footnotesize,
    breakatwhitespace=false,         
    breaklines=true,                 
    captionpos=b,                    
    keepspaces=true,                 
    numbers=left,                    
    numbersep=5pt,                  
    showspaces=false,                
    showstringspaces=false,
    showtabs=false,                  
    tabsize=2
}
\lstset{style=mystyle}


\title{Estudio Computacional de $A(n,d)$: \\ El Máximo Tamaño de un Código Binario con Longitud $n$ y Distancia Mínima $d$}
\author{Amador Jesus, Hernández Roberto, Hernández Belinda, Ramos Sarai }
\date{\today}

\begin{document}

\maketitle

\begin{abstract}
El número $A(n,d)$ representa la máxima cardinalidad posible de un código binario de longitud $n$ y distancia mínima de Hamming $d$. Este problema, central en la teoría de códigos correctores de errores, tiene aplicaciones en comunicaciones digitales y almacenamiento de datos. En este trabajo presentamos un enfoque computacional para determinar $A(n,d)$ para valores pequeños de $n$ y $d$. Reformulamos el problema como la búsqueda de una clique máxima en un grafo de Hamming, implementamos algoritmos exactos (Bron–Kerbosch) y exploramos cotas superiores mediante programación lineal. El reporte incluye el marco teórico, la metodología algorítmica, el diseño de la implementación en Python y los resultados esperados. Todo el código se encuentra en un repositorio de GitHub para facilitar la colaboración y la reproducibilidad.
\end{abstract}

\section{Introducción}
\label{sec:intro}

Un \textit{código binario} $C \subseteq \mathbb{F}_2^n$ es un conjunto de palabras (vectores) de longitud $n$ sobre el alfabeto $\{0,1\}$. La \textit{distancia de Hamming} $d_H(x,y)$ entre dos vectores $x,y\in\mathbb{F}_2^n$ es el número de coordenadas en las que difieren. La \textit{distancia mínima} del código es $d = \min_{x\neq y\in C} d_H(x,y)$. El parámetro $A(n,d)$ se define como
\[
A(n,d) \;=\; \max\{ |C| \;:\; C\subseteq\mathbb{F}_2^n,\; \forall x\neq y\in C,\; d_H(x,y)\ge d \}.
\]

Determinar $A(n,d)$ es un problema combinatorio de gran relevancia porque permite conocer el tamaño máximo de un código que puede detectar o corregir un cierto número de errores. Dado que el problema es $\mathcal{NP}$-duro (para $d$ fijo y $n$ variable), no existe una fórmula cerrada general. Sin embargo, para valores pequeños de $n$ y $d$ se han tabulado valores exactos o cotas ajustadas mediante búsqueda exhaustiva y métodos algebraicos.

En este proyecto implementamos un paquete en Python que:
\begin{enumerate}
    \item Genera el grafo de Hamming asociado a $(n,d)$,
    \item Calcula el número clique máximo (solución exacta para $n$ pequeños) mediante el algoritmo de Bron–Kerbosch,
    \item Proporciona cotas inferiores heurísticas y cotas superiores mediante programación lineal (cota de Delsarte),
    \item Valida los resultados con tablas de referencia (Brouwer, Agrell).
\end{enumerate}
El presente reporte describe los fundamentos teóricos, la metodología computacional y el plan de implementación. Todo el código fuente está versionado en GitHub, lo que permite un trabajo colaborativo eficiente entre los miembros del equipo.

\section{Marco Teórico}
\label{sec:teoria}

\subsection{Reducción a Problema de Clique en un Grafo}
\label{sub:grafo}
Definimos el \textit{grafo de Hamming} $H(n,d) = (V,E)$ donde:
\[
V = \mathbb{F}_2^n = \{0,1\}^n,\qquad 
E = \{ \{u,v\} \subset V : u\neq v,\; d_H(u,v) \ge d \}.
\]
Entonces, un código $C$ con distancia mínima $\ge d$ corresponde exactamente a un \textit{conjunto independiente} en el complemento del grafo, o equivalentemente a una \textit{clique} en $H(n,d)$. Por lo tanto,
\[
A(n,d) = \omega(H(n,d)),
\]
donde $\omega(G)$ denota el número de la clique máxima (el mayor conjunto de vértices mutuamente adyacentes). Esta reformulación permite usar algoritmos estándar de teoría de grafos.

\subsection{Cotas Fundamentales}
Para acotar $A(n,d)$ sin necesidad de una búsqueda exhaustiva, se utilizan desigualdades clásicas:

\begin{itemize}
    \item \textbf{Cota de Hamming (esfera-empaquetamiento)}:  
    \[
    A(n,d) \le \frac{2^n}{\sum_{i=0}^{\lfloor (d-1)/2\rfloor} \binom{n}{i}}.
    \]
    \item \textbf{Cota de Singleton}:  
    \[
    A(n,d) \le 2^{n-d+1}.
    \]
    \item \textbf{Cota de Plotkin}: si $d > n/2$, entonces $A(n,d) \le \left\lfloor \frac{2d}{2d-n} \right\rfloor$.
    \item \textbf{Cota de Gilbert–Varshamov (existencial)}: existe un código de tamaño al menos $\frac{2^n}{\sum_{i=0}^{d-1} \binom{n}{i}}$, por lo que $A(n,d) \ge$ esa cantidad.
\end{itemize}
Estas cotas nos proporcionan intervalos iniciales y son útiles para verificar la corrección de los resultados computacionales.

\subsection{Complejidad Computacional}
El problema de determinar $\omega(G)$ para un grafo general es $\mathcal{NP}$-duro y, bajo la hipótesis $\mathcal{P}\neq\mathcal{NP}$, no admite un algoritmo polinomial. Para $H(n,d)$, el número de vértices es $2^n$, lo que restringe los casos tratables exactamente a $n \le 10$ o $n \le 12$ (dependiendo de $d$). Para valores mayores, debemos recurrir a cotas superiores e inferiores obtenidas por métodos heurísticos o de optimización.

\section{Metodología Computacional}
\label{sec:metodo}

\subsection{Generación del Grafo de Hamming}
La generación explícita de la matriz de adyacencia requiere $O(2^{2n})$ comparaciones de distancia. Para $n\le 10$ esto es factible (poco más de $10^6$ pares). Implementamos una función que:
\begin{itemize}
    \item Enumera todos los vectores de $\mathbb{F}_2^n$ como enteros de $0$ a $2^n-1$ (representación binaria).
    \item Calcula la distancia de Hamming mediante la operación XOR y el conteo de bits (función \texttt{popcount}).
    \item Construye un grafo en memoria usando listas de adyacencia (para eficiencia).
\end{itemize}

\subsection{Algoritmos Exactos: Bron–Kerbosch}
El algoritmo de Bron–Kerbosch con pivoteo encuentra todas las cliques maximales de un grafo. Su complejidad en el peor caso es $O(3^{|V|/3})$, pero en la práctica es muy eficiente para grafos dispersos. Utilizamos una versión recursiva que mantiene tres conjuntos:
\begin{itemize}
    \item \texttt{R}: clique en construcción,
    \item \texttt{P}: candidatos que pueden extender la clique,
    \item \texttt{X}: vértices ya procesados (evita duplicados).
\end{itemize}
El pivote se elige de $P \cup X$ para minimizar las ramificaciones. El algoritmo se detiene cuando $P$ y $X$ están vacíos, registrando $R$ como una clique maximal. El tamaño máximo entre todas las cliques maximales es $\omega(G)$.

Para $n=10$, $|V|=1024$, el grafo suele ser denso cuando $d$ es pequeño, lo que puede provocar explosión combinatoria. En esos casos se aplican podas adicionales (por ejemplo, ordenar los vértices por grado degenerado).

\subsection{Enfoques Heurísticos para $n$ Moderadamente Grande}
Para $n > 12$ no es factible ejecutar Bron–Kerbosch. Proponemos dos heurísticas:
\begin{enumerate}
    \item \textbf{Búsqueda local greedy}: se parte de un vértice aleatorio y se añaden iterativamente los vértices que sean adyacentes a todos los ya seleccionados, priorizando aquellos con mayor grado.
    \item \textbf{Algoritmo de recocido simulado}: se define una función de energía negativa del tamaño de la clique y se aplican movimientos de intercambio de vértices.
\end{enumerate}
Estas heurísticas proporcionan cotas inferiores (códigos grandes) que pueden ser mejoradas con el tiempo de ejecución.

\subsection{Cotas Superiores mediante Programación Lineal}
La cota de Delsarte \cite{delsarte1973}, también conocida como cota de programación lineal, ha sido históricamente la mejor para muchos valores de $n$ y $d$. Se basa en las desigualdades de MacWilliams y en los polinomios de Krawtchouk. Resolvemos el siguiente programa lineal (dual) para obtener una cota superior:
\[
\max \; \sum_{i=0}^n \beta_i A_i \quad\text{sujeto a}\quad \beta_0=1,\; \beta_i\ge0,\; \sum_{i=0}^n \beta_i K_i(j) \ge 0\; \forall j,
\]
donde $K_i(j)$ es el polinomio de Krawtchouk y $(A_i)$ es la distribución de distancias de un código. Este programa se resuelve con bibliotecas de optimización lineal como \texttt{pulp} o \texttt{ortools}.

\section{Plan de Implementación y Resultados Esperados}
\label{sec:implementacion}

\subsection{Validación con Tablas Conocidas}
Compararemos los resultados de nuestro código con las tablas de referencia de Brouwer \cite{brouwer2018} y Agrell \cite{agrell2001}. Por ejemplo, sabemos que:
\[
A(5,3)=4,\quad A(7,3)=16,\quad A(8,3)=20.
\]
Nuestro algoritmo exacto debe reproducir estos valores para $n\le 8$. Para $n=9,10$ compararemos las cotas obtenidas por programación lineal con las reportadas en la literatura.
\documentclass{article}
\usepackage[utf8]{inputenc}
\usepackage{booktabs}
\usepackage{caption}

\begin{document}

\begin{table}[htbp]
\centering
\caption[Tabla de cotas para códigos binarios]{Tabla actualizada de cotas para códigos binarios de longitud menor a 25, con extensión para $d > 10$ \\ \small (Updated Table of Bounds for Binary Codes of Length Less than 25, with Extension for $d > 10$)}
\label{tab:binary_bounds}
\smallskip
\parbox{0.9\linewidth}{\small \textbf{Descripción:} Esta tabla muestra las mejores cotas superiores e inferiores conocidas (rangos separados por guiones) o valores exactos para códigos binarios de longitud $n$ (filas) y distancia mínima $d$ (columnas). Los valores únicos indican cotas exactas. Los rangos como $a-b$ representan el intervalo entre la cota inferior conocida $a$ y la cota superior conocida $b$. Para $n \geq 25$, se incluyen expresiones como $2^{19}-599184$ que denotan valores exactos o cotas calculadas. Esta tabla es una extensión de resultados clásicos para distancias $d>10$, actualizada hasta $n=28$.}
\medskip
\begin{tabular}{c c c c c c c c}
\toprule
$n$ & $d=4$ & $d=6$ & $d=8$ & $d=10$ & $d=12$ & $d=14$ & $d=16$ \\
\midrule
6  & 4            & 2            & 1    & 1    & 1    & 1    & 1 \\
7  & 8            & 2            & 1    & 1    & 1    & 1    & 1 \\
8  & 16           & 2            & 2    & 1    & 1    & 1    & 1 \\
9  & 20           & 4            & 2    & 1    & 1    & 1    & 1 \\
10 & 40           & 6            & 2    & 2    & 1    & 1    & 1 \\
11 & 72           & 12           & 2    & 2    & 1    & 1    & 1 \\
12 & 144          & 24           & 4    & 2    & 2    & 1    & 1 \\
13 & 256          & 32           & 4    & 2    & 2    & 1    & 1 \\
14 & 512          & 64           & 8    & 2    & 2    & 2    & 1 \\
15 & 1024         & 128          & 16   & 4    & 2    & 2    & 1 \\
16 & 2048         & 256          & 32   & 4    & 2    & 2    & 2 \\
17 & 2816-3276    & 258-340      & 36   & 6    & 2    & 2    & 2 \\
18 & 5632-6552    & 512-673      & 64   & 10   & 4    & 2    & 2 \\
19 & 10496-13104  & 1024-1237    & 128  & 20   & 4    & 2    & 2 \\
20 & 20480-26168  & 2048-2279    & 256  & 40   & 6    & 2    & 2 \\
21 & 40960-43688  & 2560-4096    & 512  & 42-47& 8    & 4    & 2 \\
22 & 81920-87333  & 4096-6941    & 1024 & 64-84& 12   & 4    & 2 \\
23 & 163840-172361& 8192-13674   & 2048 & 80-150&24   & 4    & 2 \\
24 & 327680-344308& 16384-24106  & 4096 & 136-268&48  & 6    & 4 \\
25 & $2^{19}-599184$ & 17920-47538 & 4096-5421 & 192-466 & 52-55 & 8 & 4 \\
26 & $2^{20}-1198368$ & 32768-84260 & 4104-9275 & 384-836 & 64-96 & 14 & 4 \\
27 & $2^{21}-2396736$ & 65536-157285 & 8192-17099 & 512-1585 & 128-169 & 28 & 6 \\
28 & $2^{22}-4792950$ & 131072-291269 & 16384-32151 & 1024-2817 & 178-288 & 56 & 8 \\
\bottomrule
\end{tabular}
\end{table}

\end{document}


\subsection{Resultados Esperados y Análisis}
Se espera que para $n \le 10$ y $d$ moderado ($d \ge 3$) podamos obtener el valor exacto de $A(n,d)$. Para casos más allá de la capacidad exacta, presentaremos intervalos:
\[
A_{\text{inf}}(n,d) \le A(n,d) \le A_{\text{sup}}(n,d),
\]
donde $A_{\text{inf}}$ proviene de las heurísticas y $A_{\text{sup}}$ de la cota de Delsarte. Además, mostraremos gráficas de la evolución del tamaño máximo en función de $n$ para $d$ fijo.

\section{Conclusiones y Trabajo Futuro}
\label{sec:conclusion}
En este trabajo hemos presentado una metodología computacional para abordar el problema de determinar $A(n,d)$. La reducción a problema de clique en el grafo de Hamming permite aplicar algoritmos bien estudiados, y el uso de cotas teóricas complementa los resultados exactos cuando la búsqueda exhaustiva es inviable. El paquete Python desarrollado será una herramienta didáctica y de referencia para el estudio de códigos binarios.

Como trabajo futuro, planeamos:
\begin{itemize}
    \item Incorporar la cota de programación semidefinida (SDP) descrita por Schrijver \cite{schrijver2005}, que mejora la cota de Delsarte.
    \item Paralelizar el algoritmo de Bron–Kerbosch para aprovechar múltiples núcleos.
    \item Extender el estudio a códigos no binarios (alfabetos de tamaño $q$).
\end{itemize}
Todo el código y el presente reporte están disponibles en el repositorio GitHub del equipo, favoreciendo la transparencia y la colaboración.

\begin{thebibliography}{99}

\bibitem{best1978}
M. R. Best, A. E. Brouwer, F. J. MacWilliams, A. M. Odlyzko, N. J. A. Sloane.
\newblock Bounds for binary codes of length less than 25.
\newblock \textit{IEEE Trans. Inf. Theory}, IT-24(1):81–93, 1978.

\bibitem{agrell2001}
E. Agrell, A. Vardy, K. Zeger.
\newblock A table of upper bounds for binary codes.
\newblock \textit{IEEE Trans. Inf. Theory}, 47(7):3004–3006, 2001.

\bibitem{brouwer2018}
A. E. Brouwer.
\newblock Bounds for binary codes.
\newblock Tablas en línea: \url{https://www.win.tue.nl/~aeb/codes/binary-1.html}, 2018.

\bibitem{delsarte1973}
P. Delsarte.
\newblock An algebraic approach to the association schemes of coding theory.
\newblock \textit{Philips Research Reports Supplements}, No. 10, 1973.

\bibitem{schrijver2005}
A. Schrijver.
\newblock New code upper bounds from the Terwilliger algebra and semidefinite programming.
\newblock \textit{IEEE Trans. Inf. Theory}, 51(8):2859–2866, 2005.

\bibitem{bron1973}
C. Bron, J. Kerbosch.
\newblock Algorithm 457: Finding All Cliques of an Undirected Graph.
\newblock \textit{Commun. ACM}, 16(9):575–577, 1973.

\bibitem{tomita2006}
E. Tomita, A. Tanaka, H. Takahashi.
\newblock The worst-case time complexity for generating all maximal cliques.
\newblock \textit{Theor. Comput. Sci.}, 363(3):260–267, 2006.

\bibitem{ostergard2002}
P. R. J. Östergård.
\newblock A fast algorithm for the maximum clique problem.
\newblock \textit{Discrete Applied Mathematics}, 120(1-3):197–207, 2002.

\bibitem{sagemath}
Sage Developers.
\newblock \textit{SageMath Documentation: Coding Theory}.
\newblock \url{https://doc.sagemath.org/html/en/reference/coding/}, 2024.

\end{thebibliography}

\end{document}